\documentclass[11pt]{article}
\usepackage[a4paper,margin=1in]{geometry}
\usepackage{amsmath,amssymb,amsthm,mathtools}
\usepackage{graphicx}
\usepackage[hidelinks]{hyperref}
\usepackage{microtype}
\usepackage{enumitem}

\newtheorem{theorem}{Theorem}
\newtheorem{lemma}[theorem]{Lemma}
\newtheorem{corollary}[theorem]{Corollary}
\theoremstyle{remark}
\newtheorem{remark}[theorem]{Remark}

\newcommand{\R}{\mathbb{R}}
\newcommand{\B}{\mathbb{B}}
\newcommand{\T}{\mathbb{T}}

\title{Exact Incompressible Smoothing for Unitriangular Deformations\\
and Radial Twists}
\author{Candidate partial result for arXiv:2502.01336}
\date{9 August 2026}

\begin{document}
\maketitle

\begin{center}
\fbox{\parbox{0.91\textwidth}{\textbf{Status.}
Candidate partial result; likely valid; human review requested.  The results
below solve structured subcases of Open Problem 5 in Hencl's survey, not the
general incompressible Ball--Evans approximation problem.}}
\end{center}

\section{Source question}

The source is Stanislav Hencl, \emph{Ball--Evans approximation problem:
recent progress and open problems}, arXiv:2502.01336v1 (3 February 2025)
\cite{Hencl}.  Open Problem 5 appears on PDF page 6.

\begin{center}
\includegraphics[width=\textwidth]{figures/open_problem_crop.png}
\end{center}

In transcription, the question asks whether a homeomorphism
$f\in W^{1,p}(\Omega,\R^n)$ with $J_f=1$ almost everywhere can be
approximated in the Sobolev norm by incompressible diffeomorphisms (or by
piecewise affine incompressible homeomorphisms).  The survey emphasizes that
the question is open even for $n=2$.

\section{Results and proof intuition}

Here ``incompressible'' means that the Jacobian determinant is identically
one (almost everywhere for a Sobolev map and pointwise for a smooth map).

\subsection*{Proof intuition}

There are two independent mechanisms.

First, a unitriangular deformation changes coordinate $i$ only by a function
of the later coordinates.  Its derivative is upper triangular with diagonal
entries one, so ordinary smoothing of the scalar coordinate functions cannot
change the determinant.  The same order of the coordinates gives an explicit
back-substitution inverse, so smoothing cannot destroy injectivity either.

Second, a radial twist rotates every sphere $\{|x|=r\}$ through an angle that
depends on $r$.  Smoothing the angle leaves the radius fixed.  Thus every
approximant maps the ball onto itself.  In the derivative, the only extra term
is rank one; the matrix determinant lemma reduces its contribution to a
quadratic form of a skew-symmetric matrix, hence to zero.  Freezing the angle
near the origin costs arbitrarily little in the exact weighted Sobolev energy
and makes the approximants smooth at the origin.

The first theorem covers a large triangular class and, in particular,
one-dimensional rank-one laminates.  The second theorem has the additional
features that the source and target ball are fixed and the boundary trace is
preserved.  General incompressible homeomorphisms need not have either
structure, and remain outside the argument.

\section{Unitriangular deformations}

For intervals $I_j=(a_j,b_j)$ put
\[
 Q=I_1\times\cdots\times I_n,
 \qquad Q_i=I_i\times\cdots\times I_n.
\]

\begin{lemma}[Simultaneous uniform and Sobolev smoothing]\label{lem:density}
Let $V\subset\R^d$ be a bounded open rectangle, $1\leq p<\infty$, and
$u\in W^{1,p}(V)\cap C(\overline V)$.  There are functions $u_k$, smooth on
a neighborhood of $\overline V$, such that
\[
 \|u_k-u\|_{W^{1,p}(V)}+\|u_k-u\|_{L^\infty(V)}\longrightarrow0.
\]
\end{lemma}

\begin{proof}
Reflect $u$ successively across the faces of $V$.  On a slightly larger
rectangle this gives a continuous $W^{1,p}$ extension.  Convolution with a
standard mollifier converges in $W^{1,p}$ by the usual translation
continuity argument, and converges uniformly because the continuous
extension is uniformly continuous on a compact neighborhood of
$\overline V$.
\end{proof}

\begin{theorem}[Exact smoothing of unitriangular deformations]
\label{thm:triangular}
Let $n\geq2$ and $1\leq p<\infty$.  For $i=1,\ldots,n-1$, suppose
\[
 u_i\in W^{1,p}(Q_{i+1})\cap C(\overline{Q}_{i+1}),
\]
and let $c\in\R$.  Define $F:\overline Q\to\R^n$ by
\[
 \begin{aligned}
 F_1(x)&=x_1+u_1(x_2,\ldots,x_n),\\
 F_2(x)&=x_2+u_2(x_3,\ldots,x_n),\\
 &\ \vdots\\
 F_{n-1}(x)&=x_{n-1}+u_{n-1}(x_n),\\
 F_n(x)&=x_n+c.
 \end{aligned}
\]
Then $F|_Q$ is an incompressible $W^{1,p}$ homeomorphism of $Q$ onto
$F(Q)$.  There are smooth incompressible diffeomorphisms
$F_k:Q\to F_k(Q)$ of the same unitriangular form such that
\[
 \|F_k-F\|_{W^{1,p}(Q)}+\|F_k-F\|_{L^\infty(Q)}\longrightarrow0.
\]
The conclusion holds in every dimension and for every finite $p$, including
$p=1$.
\end{theorem}

\begin{proof}
Apply Lemma~\ref{lem:density} to each $u_i$, obtaining smooth $u_{i,k}$, and
define $F_k$ by the displayed formula with $u_i$ replaced by $u_{i,k}$.
The derivative matrices $DF$ almost everywhere and $DF_k$ everywhere are
upper triangular, with every diagonal entry equal to one.  Therefore
\[
 J_F=1\quad\text{a.e.},\qquad J_{F_k}=1\quad\text{everywhere}.
\]

Injectivity does not require an estimate.  Given $y=F(x)$, first recover
$x_n=y_n-c$, then
\[
 x_{n-1}=y_{n-1}-u_{n-1}(x_n),
\]
and continue backwards until $x_1$ is recovered.  These formulas give a
continuous inverse on $F(Q)$.  The same formulas, with smooth $u_{i,k}$,
give a smooth inverse for $F_k$.  They also show directly, by the openness
of the interval constraints under the recursive recovery, that $F(Q)$ and
$F_k(Q)$ are open.  Thus $F$ is a homeomorphism onto its image and every
$F_k$ is a diffeomorphism onto its image.

Let $\pi_i(x)=(x_{i+1},\ldots,x_n)$.  Fubini's theorem gives, for either the
function or any derivative in a tail variable,
\[
 \| (u_{i,k}-u_i)\circ\pi_i\|_{L^p(Q)}^p
 =\left(\prod_{j=1}^{i}|I_j|\right)
   \|u_{i,k}-u_i\|_{L^p(Q_{i+1})}^p.
\]
The $W^{1,p}$ convergence follows after summing over $i$, while uniform
convergence is immediate from the coordinate formula.
\end{proof}

The following is the rank-one laminate form that motivated the construction.

\begin{corollary}[One-dimensional incompressible laminates]
\label{cor:laminate}
Let $U\subset\R^{n-1}$ be a bounded open rectangle and $I=(a,b)$.  Let
$B:\R^{n-1}\to\R^n$ be injective and let
$\phi\in W^{1,p}(I;\R^n)$.  Suppose
\[
 \det[Be_1\mid\cdots\mid Be_{n-1}\mid\phi'(t)]=1
 \quad\text{for a.e. }t\in I.
\]
Then
\[
 f(x',t)=Bx'+\phi(t)
\]
is an incompressible homeomorphism onto its image and is the uniform and
$W^{1,p}$ limit of smooth incompressible diffeomorphisms onto their images.
The approximating curves can be chosen to have the same endpoint values as
$\phi$.
\end{corollary}

\begin{proof}
Let
\[
 \Delta(v)=\det[Be_1\mid\cdots\mid Be_{n-1}\mid v].
\]
This is a nonzero linear functional, $\ker\Delta=\operatorname{ran}B$, and
the hypothesis says $\Delta(\phi')=1$.  Choose $v_0$ with
$\Delta(v_0)=1$.  Absolute continuity on the interval gives
\[
 \phi(t)=\phi(a)+(t-a)v_0+B w(t)
\]
for some $w\in W^{1,p}(I;\R^{n-1})$ with $w(a)=0$.  This is, after fixed
affine changes of variables, a one-layer instance of
Theorem~\ref{thm:triangular}.

For completeness, approximate $w$ in $W^{1,p}$ by smooth curves and preserve
both endpoint values by the affine correction
\[
 \widetilde w_k(t)=w_k(t)
 +\frac{b-t}{b-a}\bigl(w(a)-w_k(a)\bigr)
 +\frac{t-a}{b-a}\bigl(w(b)-w_k(b)\bigr).
\]
The trace map on a finite interval is continuous on $W^{1,p}$, so this
correction tends to zero in $W^{1,p}$.  It stays in $\R^{n-1}$, hence after
multiplication by $B$ it cannot change $\Delta(\phi_k')=1$.  Finally,
\[
 \Delta(f(x',t)-\phi(a))=t-a
\]
recovers $t$ from the image, after which injectivity of $B$ recovers $x'$.
\end{proof}

\section{Radial twists of the ball}

Let $\B^n=\{x\in\R^n:|x|<1\}$.  For $s\in\R$, let $R_s\in SO(n)$ be
rotation through angle $s$ in the $(x_1,x_2)$-plane and the identity on its
orthogonal complement.

\begin{theorem}[Fixed-ball smoothing of radial twists]
\label{thm:twist}
Let $n\geq2$, $1\leq p<\infty$, and let
$\theta:(0,1]\to\R$ be locally absolutely continuous with
\[
 \int_0^1 r^{n-1+p}|\theta'(r)|^p\,dr<\infty.
 \tag{4.1}\label{eq:weighted}
\]
Define
\[
 T_\theta(0)=0,
 \qquad T_\theta(x)=R_{\theta(|x|)}x\quad (x\ne0).
\]
Then $T_\theta$ is an incompressible $W^{1,p}$ homeomorphism of
$\B^n$ onto itself.  There exist $\theta_k\in C^\infty([0,1])$, each
constant near zero and satisfying $\theta_k(1)=\theta(1)$, such that
\[
 T_{\theta_k}\longrightarrow T_\theta
 \quad\text{in }W^{1,p}(\B^n;\R^n)\text{ and uniformly on }\overline{\B^n}.
\]
Every $T_{\theta_k}$ is a smooth incompressible diffeomorphism of $\B^n$
onto itself, and $T_{\theta_k}=T_\theta$ on $\partial\B^n$.
No limit of $\theta(r)$ as $r\downarrow0$ is required.
\end{theorem}

\begin{proof}
Write $r=|x|$, $e_r=x/r$, $Q(r)=R_{\theta(r)}$, and let a dot denote
differentiation of $R_s$ with respect to $s$.  At almost every $x\ne0$,
\[
 DT_\theta(x)
 =Q(r)+\theta'(r)\bigl(\dot R_{\theta(r)}x\bigr)\otimes e_r.
 \tag{4.2}\label{eq:derivative}
\]
Consequently $|DT_\theta|\leq C_n(1+r|\theta'|)$, and
\eqref{eq:weighted}, together with polar coordinates, gives
$DT_\theta\in L^p(\B^n)$.  Formula~\eqref{eq:derivative} is initially a
local weak-derivative formula away from the origin.  Integration by parts on
$\B^n\setminus\overline{\B(0,\varepsilon)}$ produces an inner boundary
term of size $O(\varepsilon^n)$ because $|T_\theta(x)|=\varepsilon$ there.
It vanishes as $\varepsilon\downarrow0$, proving that the displayed formula
is the global weak derivative.

The determinant is exactly one.  Since
$S_s=R_s^T\dot R_s$ is skew-symmetric, the matrix determinant lemma and
\eqref{eq:derivative} give
\[
 \begin{aligned}
 \det DT_\theta
 &=\det Q(r)\left(1+
 e_r^T Q(r)^T\theta'(r)\dot R_{\theta(r)}x\right)\\
 &=1+\frac{\theta'(r)}{r}x^TS_{\theta(r)}x=1
 \end{aligned}
\]
almost everywhere.

The identity $|T_\theta(x)|=|x|$ gives a global inverse,
\[
 T_\theta^{-1}(y)=R_{-\theta(|y|)}y.
\]
Both maps are continuous away from zero and are continuous at zero because
their norm equals the radius.  Thus $T_\theta$ is a homeomorphism of the ball
onto itself.

It remains to construct the approximation.  For $0<\delta<1/2$, freeze the
angle at radius $\delta$:
\[
 \theta^\delta(r)=
 \begin{cases}
 \theta(\delta),&0<r<\delta,\\
 \theta(r),&\delta\leq r\leq1.
 \end{cases}
\]
The maps differ only in the ball of radius $\delta$.  From
\eqref{eq:derivative}, using boundedness of rotations and their first
derivatives,
\[
 \|T_{\theta^\delta}-T_\theta\|_{W^{1,p}(\B^n)}^p
 \leq C_{n,p}\left(
 \delta^n+\int_0^\delta r^{n-1+p}|\theta'(r)|^p\,dr\right),
 \tag{4.3}\label{eq:freeze}
\]
which tends to zero.  Also
$\|T_{\theta^\delta}-T_\theta\|_\infty\leq2\delta$.

For fixed $\delta$, the function $\theta^\delta$ belongs to
$W^{1,p}(0,1)$ and is constant near zero.  Extend it constantly to the left,
mollify with radius smaller than $\delta/4$, and, if necessary, add a smooth
correction supported near $r=1$ to recover the endpoint value
$\theta(1)$.  This produces smooth $\eta_m$, constant near zero and with
$\eta_m(1)=\theta(1)$, such that
\[
 \|\eta_m-\theta^\delta\|_{W^{1,p}(0,1)}
 +\|\eta_m-\theta^\delta\|_\infty\longrightarrow0.
\]
The derivative formula gives the pointwise estimate
\[
 |DT_{\eta_m}-DT_{\theta^\delta}|
 \leq C_n\left(
 |\eta_m-\theta^\delta|+r|\eta_m'-(\theta^\delta)'|
 +r|(	heta^\delta)'|\,|\eta_m-\theta^\delta|
 \right).
\]
Hence $T_{\eta_m}\to T_{\theta^\delta}$ in $W^{1,p}$ and uniformly.
Taking a diagonal sequence with $\delta\downarrow0$ and using
\eqref{eq:freeze} yields the required $\theta_k$.

Because every $\theta_k$ is constant near zero, $T_{\theta_k}$ is smooth at
zero.  The determinant calculation is pointwise, and the explicit inverse is
smooth, so it is an incompressible diffeomorphism of the ball onto itself.
Finally $\theta_k(1)=\theta(1)$ gives equality on the boundary sphere.
\end{proof}

\begin{remark}[A low-regularity example]
The theorem permits arbitrarily many turns near the origin.  For example,
$\theta(r)=\log r$ has no limit as $r\downarrow0$, while
$r\theta'(r)=1$ and condition~\eqref{eq:weighted} holds for every finite
$p$.  The resulting homeomorphism is still uniformly and strongly Sobolev
approximable by exact incompressible diffeomorphisms of the ball.
\end{remark}

\section{Verification and scope}

The proof uses no unproved analytic or computational dependency.  The
accompanying file \texttt{verification.md} gives an adversarial step check.
The symbolic script \path{code/check_jacobians.py} verifies representative
three- and four-dimensional unitriangular examples and two- and
three-dimensional polynomial-angle radial twists.  Those computations are
sanity checks only; the proof is the triangular and rank-one determinant
calculation above.

The triangular result does not generally keep the image $F(Q)$ fixed while
smoothing; its approximants are diffeomorphisms onto their own images.  The
radial-twist result does keep both domain and image equal to $\B^n$, and it
preserves the boundary trace.  Neither result handles an arbitrary
incompressible Sobolev homeomorphism, so Open Problem 5 remains open in
general.

\subsection*{Bounded novelty check}

On 9 August 2026 we searched the run's registry, solution, attempt, and
proof-gap indexes for arXiv:2502.01336 and for the terms ``incompressible'',
``volume preserving'', ``triangular'', ``laminate'', ``shear'', and ``radial
twist''.  We also searched arXiv/web indexes for Sobolev approximation by
volume-preserving diffeomorphisms and exact determinant-one smoothing of
shears or laminates.  No exact match for either theorem above was found.

Two nearby results delimit the claim.  Mueller proves uniform approximation
results for volume-preserving homeomorphisms in a high-dimensional manifold
setting, but the cited abstract does not assert strong $W^{1,p}$ convergence
\cite{Mueller}.  Campbell, D'Onofrio, and Vitek prove strong Sobolev
diffeomorphic approximation of general locally finite piecewise affine
homeomorphisms in dimensions three and four, including simultaneous inverse
control, but their theorem does not assert exact incompressibility
\cite{CDV}.  The present result is therefore recorded only as a candidate new
structured partial theorem.  The search was bounded rather than exhaustive,
and the mechanism may be folklore.

\subsection*{Human-review recommendation}

Review as a likely valid partial result.  The most useful checks are:
\begin{enumerate}[leftmargin=2em]
\item confirm the simultaneous uniform/$W^{1,p}$ density lemma at $p=1$;
\item verify the distributional derivative across the origin in
Theorem~\ref{thm:twist};
\item verify that the fixed-ball and boundary assertions use only radius
preservation and the endpoint condition;
\item assess whether exact determinant-preserving triangular or radial-twist
smoothing is already recorded in the specialist literature.
\end{enumerate}

\begin{thebibliography}{9}

\bibitem{Hencl}
S. Hencl,
\emph{Ball--Evans approximation problem: recent progress and open problems},
arXiv:2502.01336v1, 2025.

\bibitem{Mueller}
S. Mueller,
\emph{Uniform approximation of homeomorphisms by diffeomorphisms},
arXiv:0901.1002, 2009.

\bibitem{CDV}
D. Campbell, L. D'Onofrio, and T. Vitek,
\emph{Diffeomorphic approximation of piecewise affine homeomorphisms},
arXiv:2507.02854, 2025; Journal of Geometric Analysis 36 (2026), Article 157.

\end{thebibliography}

\end{document}
